\begin{frame}{Βοηθητικοί πυρήνες: \en{Keccak} \& αποκωδικοποίηση}
  % src: results/reports/summaries/speedups_simulation.md (Tables 4-5)
  \begin{columns}[T]
    \begin{column}{0.5\textwidth}
      \textbf{\en{Keccak}/SHA-3}
      \begin{itemize}
        \item 2-way \en{SoA} με \Zvbb (\vror).
        \item \pos{\speedup{1.36}}–\pos{\speedup{1.37}}.
      \end{itemize}
    \end{column}
    \begin{column}{0.5\textwidth}
      \textbf{Αποκωδικοποίηση \en{RM}+\en{RS}}
      \begin{itemize}
        \item \RVV + \Zvbc $\mathrm{GF}(2^8)$.
        \item \pos{\speedup{13.2}}–\pos{\speedup{15.9}}.
      \end{itemize}
    \end{column}
  \end{columns}
  \vspace{0.4em}
  \begin{center}\footnotesize Μετά τον \en{polymul}, οι πυρήνες αυτοί κυριαρχούν στο εναπομείναν κόστος.\end{center}
\end{frame}
